\section{Exploración multiescala y generalización entre tareas}
\subsection{Estrategia de exploración en varias etapas}
La Slide 07 presenta el multi-stage exploration framework: divide la exploration policy en dos niveles, coarse-grained y fine-grained, que controlan respectivamente los objetivos y las secuencias de acciones de distinta escala. Esta estructura jerárquica permite descubrir macro-level goal cluster y, a la vez, converger automáticamente a los pasos de ejecución micro-level.

\begin{figure}[H]
\centering
\includegraphics[width=0.85\textwidth]{slides-images/slide-07.jpg}
\caption{La Slide 7 explica cómo la multi-stage exploration amplía de forma coordinada el goal space.}
\end{figure}

\begin{knowledgebox}{Multi-Stage Exploration}
\begin{itemize}
\item La coarse policy muestrea goal cluster no vistos;\\
\item La fine policy se encarga de las acciones detalladas dentro de ese cluster;\\
\item Se actualizan periódicamente los cluster centroids based on success rate para mantener la coverage.
\end{itemize}
\end{knowledgebox}

\subsection{Métricas de generalización entre tareas}
La Slide 08 propone una transfer matrix para evaluar la capability de la policy en distintas combinaciones de goal; las métricas principales son success coverage, policy entropy y transfer ratio. El ponente recomienda mostrar en los informes experimentales la performance tanto en los base goals como en los transfer goals, para evitar el sobreajuste a un solo escenario.

\begin{tabularx}{\textwidth}{lX}
\toprule
Métrica & Descripción \\
\midrule
Success coverage & Proporción de stored goals en los que se alcanza el éxito; un valor inferior a 70\% indica una coverage gap; \\
Policy entropy & Supervisa la diversidad de la policy en distintos goals; un valor demasiado bajo indica que es degenerate; \\
Transfer ratio & La improvement que la reward de un goal produce en otro goal; mide la generalización; \\
\bottomrule
\end{tabularx}

\subsection{Resumen de la sección}
La exploración multiescala y las métricas de generalización ayudan al equipo a identificar en qué goals hay que romper cuellos de botella estrechos y ofrecen herramientas para medir la eficacia de la transfer.
\newpage
